\textbf{Phần 1. Trắc nghiệm lựa chọn}
\setcounterpageref{ex}{0}

\begin{ex} %[5D2N11]
Khi so sánh hai số thập phân $3{,}5$ và $2{,}75$, ta thấy
\choice
{$3{,}5 < 2{,}75$ vì $3{,}5$ có ít chữ số hơn}
{$3{,}5 = 2{,}75$}
{$3{,}5 < 2{,}75$ vì $5 < 75$}
{\True $3{,}5 > 2{,}75$ vì ở phần nguyên có $3 > 2$}
\loigiai{
So sánh phần nguyên của hai số: $3 > 2$ nên $3{,}5 > 2{,}75$.\\
\textbf{Đáp án của câu hỏi là D}
}
\end{ex}

\begin{ex} %[5D2N11]
Khi so sánh hai số thập phân $2{,}75$ và $2{,}29$, ta thấy
\choice
{$2{,}75 < 2{,}29$}
{\True $2{,}75 > 2{,}29$ vì ở hàng phần mười có $7 > 2$}
{$2{,}75 = 2{,}29$}
{$2{,}75 < 2{,}29$ vì ở hàng phần trăm có $5 < 9$}
\loigiai{
Hai số có cùng phần nguyên là $2$. So sánh hàng phần mười: $7 > 2$ nên $2{,}75 > 2{,}29$.\\
\textbf{Đáp án của câu hỏi là B}
}
\end{ex}

\begin{ex} %[5D2H11]
Số thập phân nào dưới đây bằng với số $0{,}7$?
\choice
{$0{,}07$}
{$7{,}0$}
{$0{,}007$}
{\True $0{,}70$}
\loigiai{
Nếu viết thêm chữ số $0$ vào tận cùng bên phải phần thập phân thì được một số thập phân bằng nó, do đó $0{,}7 = 0{,}70$.\\
\textbf{Đáp án của câu hỏi là D}
}
\end{ex}

\begin{ex} %[5D2H11]
Số thập phân $8{,}6100$ khi bỏ các chữ số $0$ ở tận cùng bên phải phần thập phân thì được số gọn hơn là
\choice
{$8{,}6$}
{$86{,}1$}
{\True $8{,}61$}
{$8{,}610$}
\loigiai{
Bỏ hai chữ số $0$ ở tận cùng bên phải phần thập phân ta được $8{,}6100 = 8{,}61$.\\
\textbf{Đáp án của câu hỏi là C}
}
\end{ex}

\begin{ex} %[5D2H11]
Cặp số thập phân nào sau đây là \textbf{bằng nhau}?
\choice
{$0{,}4$ và $0{,}04$}
{\True $21{,}4050$ và $21{,}405$}
{$5{,}82$ và $5{,}802$}
{$13{,}05$ và $13{,}50$}
\loigiai{
Vì bỏ chữ số $0$ ở tận cùng bên phải phần thập phân của $21{,}4050$ ta được $21{,}405$.\\
\textbf{Đáp án của câu hỏi là B}
}
\end{ex}

\begin{ex} %[5D2H11]
Trong các số thập phân $3{,}604;\; 2{,}875;\; 2{,}857;\; 3{,}106$, số bé nhất là
\choice
{$3{,}604$}
{$2{,}875$}
{\True $2{,}857$}
{$3{,}106$}
\loigiai{
So sánh phần nguyên: $2 < 3$. So sánh hai số có phần nguyên là $2$: $2{,}857 < 2{,}875$ (vì $5 < 7$ ở hàng phần trăm). Vậy số bé nhất là $2{,}857$.\\
\textbf{Đáp án của câu hỏi là C}
}
\end{ex}

\begin{ex} %[5D2H11]
Sắp xếp các số $3{,}604;\; 2{,}875;\; 2{,}857;\; 3{,}106$ theo thứ tự từ bé đến lớn là
\choice
{\True $2{,}857;\; 2{,}875;\; 3{,}106;\; 3{,}604$}
{$2{,}875;\; 2{,}857;\; 3{,}106;\; 3{,}604$}
{$3{,}604;\; 3{,}106;\; 2{,}875;\; 2{,}857$}
{$2{,}857;\; 2{,}875;\; 3{,}604;\; 3{,}106$}
\loigiai{
Ta có $2{,}857 < 2{,}875 < 3{,}106 < 3{,}604$.\\
\textbf{Đáp án của câu hỏi là A}
}
\end{ex}

\begin{ex} %[5D2H11]
Số thập phân nào sau đây lớn hơn $1{,}036$ và bé hơn $2$?
\choice
{$0{,}95$}
{$1{,}036$}
{$2{,}01$}
{\True $1{,}36$}
\loigiai{
So sánh: $0{,}95 < 1{,}036 < 1{,}36 < 2 < 2{,}01$. Vậy số thỏa mãn là $1{,}36$.\\
\textbf{Đáp án của câu hỏi là D}
}
\end{ex}

\begin{ex} %[5D2H11]
Điền dấu thích hợp vào chỗ chấm: $135{,}74 \dots 135{,}73$
\choice
{\True $>$}
{$<$}
{$=$}
{Không so sánh được}
\loigiai{
Hai số có cùng phần nguyên $135$ và chữ số hàng phần mười là $7$. So sánh hàng phần trăm: $4 > 3 \Rightarrow 135{,}74 > 135{,}73$.\\
\textbf{Đáp án của câu hỏi là A}
}
\end{ex}

\begin{ex} %[5D2H11]
Chữ số thích hợp điền vào dấu hỏi chấm để $13{,}7 = 13{,}7?$ là
\choice
{$1$}
{$7$}
{\True $0$}
{$3$}
\loigiai{
Để $13{,}7 = 13{,}7?$ thì chữ số tận cùng bên phải phải là $0$ ($13{,}7 = 13{,}70$).\\
\textbf{Đáp án của câu hỏi là C}
}
\end{ex}

\begin{ex} %[5D2V11]
Ba bạn Mị, Núi, Páo đi khám sức khoẻ. Cân nặng của ba bạn lần lượt là một trong các số: $31{,}9\text{ kg};\; 32{,}5\text{ kg};\; 34{,}7\text{ kg}$. Biết bạn Núi nặng nhất, bạn Páo nhẹ nhất. Cân nặng của bạn Mị là
\choice
{\True $32{,}5\text{ kg}$}
{$31{,}9\text{ kg}$}
{$34{,}7\text{ kg}$}
{$33{,}0\text{ kg}$}
\loigiai{
So sánh cân nặng: $31{,}9\text{ kg} < 32{,}5\text{ kg} < 34{,}7\text{ kg}$.\\
Vì Núi nặng nhất ($34{,}7\text{ kg}$) và Páo nhẹ nhất ($31{,}9\text{ kg}$) nên bạn Mị nặng $32{,}5\text{ kg}$.\\
\textbf{Đáp án của câu hỏi là A}
}
\end{ex}

\begin{ex} %[5D2V11]
Chiều dài của ba cây cầu: Cầu Long Biên ($2{,}29\text{ km}$), Cầu Cần Thơ ($2{,}75\text{ km}$), Cầu An Đông ($3{,}5\text{ km}$). Cây cầu dài nhất là
\choice
{Cầu Long Biên}
{\True Cầu An Đông}
{Cầu Cần Thơ}
{Cả ba cây cầu dài bằng nhau}
\loigiai{
So sánh: $2{,}29 < 2{,}75 < 3{,}5$. Cây cầu dài nhất là Cầu An Đông ($3{,}5\text{ km}$).\\
\textbf{Đáp án của câu hỏi là B}
}
\end{ex}

\textbf{Phần 2. Trắc nghiệm đúng sai}
\setcounterpageref{ex}{0}

\begin{ex} %[5D2H11]
Cho các số thập phân: $x = 8{,}9$; $y = 5{,}82$; $z = 0{,}17$; $t = 31{,}6$. Xét tính đúng sai của các khẳng định sau:
\choiceTF
{\True Viết số $x$ dưới dạng phần thập phân có ba chữ số ta được $8{,}900$.}
{Viết số $y$ dưới dạng phần thập phân có ba chữ số ta được $5{,}082$.}
{\True Số lớn nhất trong bốn số đã cho là $t = 31{,}6$.}
{\True Số bé nhất trong bốn số đã cho là $z = 0{,}17$.}
\loigiai{
\begin{itemize}
    \item Thêm hai chữ số $0$ vào bên phải phần thập phân của $8{,}9$ ta được $8{,}900$.
    \item Thêm một chữ số $0$ vào tận cùng bên phải của $5{,}82$ ta được $5{,}820 \neq 5{,}082$.
    \item So sánh: $0{,}17 < 5{,}82 < 8{,}9 < 31{,}6 \Rightarrow$ Số lớn nhất là $31{,}6$, số bé nhất là $0{,}17$.
\end{itemize}
\begin{itemize}
    \item a) \textbf{Đúng}.
    \item b) \textbf{Sai}: Kết quả là $5{,}820$.
    \item c) \textbf{Đúng}.
    \item d) \textbf{Đúng}.
\end{itemize}
}
\end{ex}

\begin{ex} %[5D2V11]
Ba bạn Mị, Núi, Páo có cân nặng lần lượt là $32{,}5\text{ kg}; 34{,}7\text{ kg}; 31{,}9\text{ kg}$. Xét tính đúng sai của các khẳng định sau:
\choiceTF
{\True Bạn Núi nặng hơn bạn Mị là $2{,}2\text{ kg}$.}
{Tổng cân nặng của bạn Páo và bạn Mị là $64{,}5\text{ kg}$.}
{\True Bạn Páo nhẹ hơn bạn Núi là $2{,}8\text{ kg}$.}
{\True Trung bình mỗi bạn nặng $33{,}033\dots\text{ kg}$ (lớn hơn $33\text{ kg}$).}
\loigiai{
\begin{itemize}
    \item Núi nặng hơn Mị: $34{,}7 - 32{,}5 = 2{,}2\text{ kg}$.
    \item Tổng cân nặng Páo và Mị: $31{,}9 + 32{,}5 = 64{,}4\text{ kg} \neq 64{,}5\text{ kg}$.
    \item Páo nhẹ hơn Núi: $34{,}7 - 31{,}9 = 2{,}8\text{ kg}$.
    \item Tổng cân nặng $3$ bạn: $31{,}9 + 32{,}5 + 34{,}7 = 99{,}1\text{ kg}$. Trung bình mỗi bạn: $99{,}1 : 3 \approx 33{,}03\text{ kg} > 33\text{ kg}$.
\end{itemize}
\begin{itemize}
    \item a) \textbf{Đúng}.
    \item b) \textbf{Sai}: Tổng là $64{,}4\text{ kg}$.
    \item c) \textbf{Đúng}.
    \item d) \textbf{Đúng}.
\end{itemize}
}
\end{ex}

\textbf{Phần 3. Trả lời ngắn}
\setcounterpageref{ex}{0}

\begin{ex} %[5D2H11]
Tìm chữ số $x$ thích hợp để $8{,}6100 = 8{,}6x$.
\shortans{1}
\loigiai{
\textbf{Đáp số:} $1$\\
Bỏ các chữ số $0$ tận cùng bên phải phần thập phân ta có $8{,}6100 = 8{,}61 \Rightarrow x = 1$.
}
\end{ex}

\begin{ex} %[5D2H11]
Có bao nhiêu số thập phân có một chữ số ở phần thập phân lớn hơn $3{,}5$ và nhỏ hơn $4{,}2$?
\shortans{6}
\loigiai{
\textbf{Đáp số:} $6$\\
Các số đó là: $3{,}6;\; 3{,}7;\; 3{,}8;\; 3{,}9;\; 4{,}0;\; 4{,}1$. Có tất cả $6$ số.
}
\end{ex}

\begin{ex} %[5D2H11]
Tìm số tự nhiên $x$ thỏa mãn: $8{,}9 < x < 9{,}8$.
\shortans{9}
\loigiai{
\textbf{Đáp số:} $9$\\
Số tự nhiên duy nhất nằm giữa $8{,}9$ và $9{,}8$ là $9$.
}
\end{ex}

\begin{ex} %[5D2H11]
Tìm số tự nhiên $n$ lớn nhất thỏa mãn: $n < 5{,}82$.
\shortans{5}
\loigiai{
\textbf{Đáp số:} $5$\\
Vì $n$ là số tự nhiên và $n < 5{,}82$ nên giá trị lớn nhất của $n$ là $5$.
}
\end{ex}

\begin{ex} %[5D2V11]
Cầu Cần Thơ dài $2{,}75\text{ km}$, Cầu Long Biên dài $2{,}29\text{ km}$. Hỏi cầu Cần Thơ dài hơn cầu Long Biên bao nhiêu mét?
\shortans{460}
\loigiai{
\textbf{Đáp số:} $460$\\
Đổi $2{,}75\text{ km} = 2\,750\text{ m}$ và $2{,}29\text{ km} = 2\,290\text{ m}$.\\
Hiệu chiều dài là: $2\,750 - 2\,290 = 460\text{ (m)}$.
}
\end{ex}

\begin{ex} %[5D2V11]
Tổng cân nặng của ba bạn Mị, Núi và Páo là bao nhiêu ki-lô-gam biết cân nặng ba bạn lần lượt là $32{,}5\text{ kg}; 34{,}7\text{ kg}; 31{,}9\text{ kg}$? (Nhập kết quả dưới dạng số thập phân).
\shortans{99,1}
\loigiai{
\textbf{Đáp số:} $99{,}1$\\
Tổng cân nặng là: $32{,}5 + 34{,}7 + 31{,}9 = 99{,}1\text{ (kg)}$.
}
\end{ex}
